\section{A Morse-Theoretic Filtration}
\subsection{Conley index}
\begin{definition}[Compact Invariant Isolated Subset] \label{def:morseFiltr-compactInvariantIsolatedSubset}
	For a flow $\phi_t:M\to M$ on a locally compact metric space, we call a subset $S\subseteq M$ an \textbf{invariant subset} if and only if $\phi_t(S)=S$ for all $t\in \R$. For any subset $N\subseteq M$ we define the \textbf{maximal invariant subset}
	\begin{align*}
		I(N)&=\{x\in N \mid \phi_t(x)\in N \forall t\in \R\}\\
		&=\bigcap_{t\in \R} \phi_t(N).
	\end{align*} A \textbf{compact invariant subset} $S$ is called \textbf{isolated} if a compact neighbourhood $N$ exists such that $I(N)=S$.
\end{definition}
\begin{definition}[Index Pairs] \label{def:morseFiltr-indexPairs}
	Let $S$ be an isolated compact invariant subset. A topological pair $(N,L)$ of compact subsets of $M$ where $L\subseteq N$ is called an \textbf{index pair of }S, if it satisfies the following:
	\begin{enumerate}
		\item $S=I(\overline{N\setminus L})\subseteq (\mathring{N\setminus L})$.
		\item  $x\in L$ and $\phi_{[0,t]}(x)\subseteq N$ implies that  $\phi_{[0,t]}(x)\subseteq L$. We call $L$ \textbf{positively invariant in } $N$. Here $\phi_{[0,t]}(x)\coloneq\{ \phi_{\tilde{t}}(x)|\tilde{t}\in [0,t]\}$.
		\item For all $x\in N$ such that a $t$ exists, with $\phi_t(x)\not\in N$, there exists a $t'$ with $\phi_{[0,t']}(x)\subseteq N$ and $\phi_{t'}(x)\in L$.
	\end{enumerate}
	This definition captures how the flow lines leave the invariant space $S$. Due to the third property we call $L$ the \textbf{exit set}.
\end{definition}


\begin{definition}[Regular Index Pairs]\label{def:morseFiltr-regularIndexPairs}
	We call an index pair $(N,L)$ \textbf{regular} if and only if the inclusion $i: L \hookrightarrow N$ is a cofibration.
\end{definition}
\begin{theorem}[Homotopy Equivalence of Index Pairs]\label{thm:morseFiltr-homotopyEquivalenceOfIndexPairs}
	If $S$ is an isolated compact invariant subset and $(N,L)$ and $(\tilde{N},\tilde{L})$ are two index pairs of $S$, then $N/L$ and $\tilde{N}/\tilde{L}$ are homotopy equivalent as pointed spaces. This homotopy equivalence is canonical, and transitive up to homotopy. 
\end{theorem}
\begin{proof}
	For a proof of the homotopy equivalence we refer to \cite{banyaga2004lectures}, Theorem 7.15. The constructed homotopy equivalence is given as follows: Choose $T\geq 0$ such that the following implications hold for $t\geq T$:
\begin{align*}
\phi_{[-t,t]}(x) \subseteq N\setminus L \quad &\Rightarrow \quad x \in \tilde{N}\setminus \tilde{L}\,, \\
\phi_{[-t,t]}(x) \subseteq  \tilde{N}\setminus \tilde{L} \quad &\Rightarrow \quad x \in N\setminus L\,, 
\end{align*} where $\phi_{[a,b]}(x) \coloneq \{\phi_t(x)\mid t\in [a,b]\}$.
Then the map:
\begin{align*}
h:N/L \times [T,\infty) & \to \tilde{N}/\tilde{L} \\
([x],t) & \mapsto 
\begin{cases}
[\phi_{3t}(x)] & \text{ if } \phi_{[0,2t]}(x)\subseteq N\setminus L \text{ and } \phi_{[t,3t]}(x)\subseteq \tilde{N}\setminus \tilde{L}\\
[\tilde{L}] & \text{ otherwise,}
\end{cases}
\end{align*}
is continuous and for any fixed $t$ this gives our homotopy equivalence. Obviously, such a map is canonical and transitive up to homotopy.
\end{proof}

\begin{definition}[Connected Simple System]\label{def:morseFiltr-connectedSimpleSystem}
	A \textbf{connected simple system} consists of a set $I_0$ of pointed spaces along with a set $I_m$ of homotopy classes of maps between these spaces such that 
		\begin{enumerate}[(i)]
			\item $\hom(X,Y)\coloneq \{[f]\in [X,Y]\mid[f]\in I_m\}$ consists of exactly one element for each pair of spaces $(X,Y)\in I_0^2$, where $[X,Y]$ denotes the set of homotopy classes of maps between $X$ and $Y$.
			\item If $X,Y,Z\in I_0,[f]\in \hom(X,Y),$ and $[g]\in \hom(Y,Z)$, then $[g\circ f]\in \hom(X,Z)$.
			\item $\hom(X,X)=\{[\id_X]\}$ for all $X\in I_0$. 
		\end{enumerate}
\end{definition}
\begin{remark}\label{rem:morseFiltr-connectedSimpleSystemSalamon}
	This definition comes from Salamon in \cite{Salamon1985}.
\end{remark}


\begin{remark}\label{rem:morseFiltr-indexPairConnectedSimpleSystem} Let $S$ be an isolated compact invariant subset of $M$ with respect to the Morse flow $\phi$. Then the collections
\begin{align*}
	I_0&\coloneq \{N  / L\mid (N,L)\text{ is an index pair of }S\}\\
	I_m&\coloneq \{[f] \mid f \text{ is a homotopy equivalence from Theorem \ref{thm:morseFiltr-homotopyEquivalenceOfIndexPairs}}\}\,,
\end{align*} make up a connected simple system. We denote this by $\mathcal{IP}_S$.
\end{remark}


\begin{definition}[Morphisms of Connected Simple Systems]\label{def:morseFiltr-morphismsOfConnectedSimpleSystems}
A \textbf{morphism} $\Phi:I\to J$ between two connected simple systems $I=(I_0,I_m)$ and $J=(J_0,J_m)$ is a set of homotopy classes of maps between spaces in $I_0$ and $J_0$ such that
\begin{enumerate}[(i)]
	\item For every $X\in I_0$ and $Y\in J_0$ the set $\Phi(X,Y)\coloneq \{[\phi]\in [X,Y]\mid[\phi]\in \Phi\}$ consists of exactly one element. 
	\item If $X,\bar{X}\in I_0$ and $Y,\bar{Y}\in J_0$ and if $[\phi]\in \Phi(X,Y)$, $[f]\in \hom(\bar{X},X)$, $[g]\in \hom(Y,\bar{Y})$, then $[g\circ \phi\circ f]\in \Phi(\bar{X},\bar{Y})$
\end{enumerate}
\end{definition}
\begin{remark}\label{rem:morseFiltr-morphismInducedBySingleMap}
	By the second property of a morphism between connected simple systems $I$ and $J$, we can induce such a morphism by a single map $\phi:X\to Y$ where $X\in I_0$ and $Y\in J_0$. 
\end{remark}


\begin{definition}[Canonical Isomorphism Class]\label{def:morseFiltr-canonicalIsomorphismClass}
	A \textbf{canonical isomorphism class} of Abelian groups (or other algebraic structures) is a set of groups $R_0$ together with a set of group homomorphisms $R_m$ such that
	\begin{enumerate}[(i)]
		\item For $R,L \in R_0$ the set $\hom(R,L)\coloneq \{f:R\to L: f \in R_m\}$ consists of exactly one element. 
		\item If $R,L,S\in R_0$ and $f\in \hom(R,L)$ and $g\in \hom(L,S)$, then $g\circ f\in \hom(L,S)\,$.
		\item $\hom(R,R)=\{\id_R\}$
	\end{enumerate}
	A \textbf{morphism} between two canonical isomorphism classes of Abelian groups $\Psi: (R_0,R_m)\to (S_0,S_m)$ is a collection of homomorphisms such that 
	\begin{enumerate}[(i)]
		\item For any $R\in R_0$ and $S\in S_0$ the set $\Psi(R,S)\coloneq \{f:R\to S: f\in \Psi\}$ contains exactly one element.
		\item If $R,\bar{R}\in R_0$ and $S,\bar{S}\in S_0$, $f\in \hom(\bar{R},R)$, $g\in \hom(S,\bar{S})$ and $\psi\in \Psi(R,S)$, then $g\circ \psi\circ f\in \Psi(\bar{R},\bar{S})\,$.
	\end{enumerate}
\end{definition}

\begin{definition}\label{def:morseFiltr-systemCategories}
We denote the category of connected simple systems with their morphisms by $\CSS$ and similar we call the category of canonical isomorphism classes of Abelian groups by $\CIAB$, respectively, the category for rings by $\CIRing$.
\end{definition}

\begin{definition} \label{def:morseFiltr-postcomposedFunctor}
	Let $G$ be an Abelian group and $X$ be any set. We can \textbf{tag} $G$ with $X$ by defining the set $G\times \{X\}$ and endowing it with the algebraic information from $G$. (for example we define $(g,X)\cdot (g',X)=(g\cdot_Gg',X)$). Let $f:G\to L$ be a homomorphism. Then $f$ induces a homomorphism between the tagged $G\times \{X\}$ and $L\times \{Y\}$ by sending $(g,X)\mapsto (f(g),Y)$.
	Let $F:\TopP\to \Ab$ be a (contravariant) functor from pointed topological spaces to Abelian groups that factors over the homotopy category, i.e. is homotopy invariant. We first define the tagged functor $F^{tag}$ to map spaces to groups that are tagged with the spaces they come from:
	\begin{align*}
		F^{tag}: \TopP &\to \Ab\\
		X&\mapsto F(X)\times \{X\}\\
		\big(f:X\to Y\big)&\mapsto \big(F^{tag}(f):F(Y)\times \{Y\}\to F(X)\times\{X\}\big)
	\end{align*}
	For covariant functors the procedure is analogous.
	Now we can define the \textbf{postcomposed functor} $F_{\#}$ from the category of connected simple systems to the canonical isomorphism classes as follows:
	\begin{align*}
	F_{\#} : \CSS		&\to   \CIAB \\
	I_0 & \mapsto \{F^{tag}(X)\mid X \in I_0 \} \\
	I_m & \mapsto \{ F^{tag}(g) \mid [g]\in I_m \}\\
	\end{align*}
	Let $\Phi:(I_0,I_m)\to (J_0,J_m)$ be a morphism in $\CSS$. Then we define
	\begin{equation*}
		F_{\#}(\Phi): F_{\#}(J_0,J_m) \to  F_{\#}(I_0,I_m)\,,\quad 
		F_{\#}(\Phi) \coloneq \{F^{tag}(g)\mid [g]\in \Phi \}\,.
	\end{equation*}
This obviously becomes a new functor, where the tagging removed uniqueness problems for non-injective (in the literal sense not up to isomorphism) functors $F$. If $F(X)=F(Y)$ for different $X,Y\in I_0$ with $g:X\to Y\in I_m$, then without tagging there could be two different morphisms $F(X)\to F(X)$: the identity and $F(g)$. Since our functors are almost always injective, we will ignore the tagging from now on. 
\end{definition}
\begin{remark}
	Obviously, the procedure of postcomposing a functor works for any homotopy invariant functor from pointed topological spaces into any algebraic category and not just Abelian groups.
\end{remark}

\begin{remark}[Critical points as isolated compact invariant subsets] \label{rem:morseFiltr-criticalPointsAsIsolatedCompactInvariantSubsets}
	Let ${f:M\to \R}$ be a Morse function on a smooth Riemannian manifold $(M,g)$. Let $q\in \Crit_k(f)$. Around $q$ we have refined Morse charts $\phi:U \to T_qM$ (after identifying $x_i$ and $\partial_i$) where  ${\phi(W(q\to)\cap U)\subseteq T_q^{u}M}$ and ${\phi(W(\to q)\cap U)\subseteq T_q^{s}M}$. With this we can define the balls:
	\begin{align*}
		D^{s}_\epsilon &~\coloneq~\{v\in T_q^{s}M \mid \norm{v}\leq \epsilon \}, \\
		D^{u}_\epsilon &~\coloneq~\{v\in T_q^{u }M \mid \norm{v}\leq \epsilon \}. 
	\end{align*}
For epsilon small enough, they give rise to the index pair $N_q\coloneq\phi^{-1}( D^{s }_\epsilon\times  D^{u }_\epsilon)$ and $L_q\coloneq\phi^{-1}( D^{ s}_\epsilon\times \partial D^{u}_\epsilon)$. To see this, notice that in such a refined Morse chart $\phi$, the differential of the flow in the critical point splits the tangent space into $T^u_qM\times T^s_qM$. Finally, since the property of $\diff \phi_x$ being a contraction on $T^s_qM$ is an open condition and thereby there is an $\epsilon>0$ such that $\diff \phi_x$ is a contraction on $T^s_qM$ for all $x$ with $\norm{\phi^{-1}(x)}<\epsilon$. Thereby, the norm of $\phi_t(x)$ projected to $T^s_qM$ is strictly decreasing in $t$. This proves that $L_q$ is an exit set of $N_q$. Similarly, one sees that $L_q$ is positively invariant in $N_q$. 

This index pair is in fact regular, since $(N_q,L_q)\cong  (D^{ s}_\epsilon\times  D^{u }_\epsilon, D^{s}_\epsilon\times \partial D^{u}_\epsilon) \cong (D^{u}_{\epsilon},\partial D^{u}_{\epsilon} )$.\todo{ warum ist das eine cofibration? } Now we know the K-group of such a pair:
	\begin{align} \label{eq: homology of index pair of critical point}
		K^i(N_q,L_q)\cong K^i(D^{k},\partial D^{k})\cong \begin{cases}
			\Z & \text{if } i-k\in 2\Z,\\
			0& \text{else. }
		\end{cases}
	\end{align}
	Notice that for $k=0$ we need to consider $\partial D^k=\emptyset$ to get an index pair. That's why we let $\partial$ denote the manifold boundary instead of the topological boundary. However, now we can identify for all $k$: \begin{equation} \label{ kritical points as direct sum of homology groups}
		C^{k}(M,\mathfrak{A},\tilde{K}^{l}(pt))=\bigoplus_{ \Crit_k(f)} \Z \cong \bigoplus_{q\in \Crit_k(f)} K^k(N_q,L_q).
	\end{equation}
	Or using the language of connected simple systems, we have the identification 
	\begin{equation} 
		C^{k}(M,\mathfrak{A},\tilde{K}^{l}(pt))\cong \bigoplus_{q\in \Crit_k(f)} K_{\#}^k(\mathcal{IP}_q)\,,
	\end{equation}where the isomorphism can be read as an isomorphism of canonical isomorphism classes with the left-hand side vierwd as an object of $\CIAB$ consists of only one group.
\end{remark}